\documentclass[a4paper,12pt]{article}
\usepackage[utf8]{inputenc}
\usepackage[T1]{fontenc}
\usepackage[french]{babel}
\usepackage{geometry}
\usepackage{amsmath}
\usepackage{amssymb}
\usepackage{array}
\usepackage{graphicx}
\usepackage{booktabs}
\usepackage{pgfplotstable}
\usepackage{float}
\usepackage{pgfplots}
\pgfplotsset{compat=1.18}

% Adjust margins to fit tables
\geometry{hmargin=2cm, vmargin=2cm}

\title{\textbf{Compte Rendu : Pendule de Torsion}}
\author{Ilyas Zanam \\ Noha Maghfoul}
\date{}

\begin{document}

\begin{titlepage}
    \centering
    \vspace*{2cm}
    
    {\Huge\bfseries Compte rendu Tp N°1\par}
    \vspace{1cm}
    {\LARGE Pendule de Torsion\par}
    \vspace{2cm}
    
    % Add your image here
    % \includegraphics[width=0.4\textwidth]{pnd.png}
    
    \vspace{2cm}
    
    {\Large réalisé par:\par}
    \vspace{0.5cm}
    {\large Ilyas Zanan \hspace{2cm} Noha Maghfoul\par}
    
    \vspace{1.5cm}
    
    {\Large encadré par:\par}
    \vspace{0.5cm}
    {\large Pr. Mohamed AHD\par}
    
    \vspace{1.5cm}
    
    {\Large responsable:\par}
    \vspace{0.5cm}
    {\large Pr. Mohamed GOUIGHRI\par}
    
    \vfill
    
\end{titlepage}

\newpage

\begin{center}
    \textbf{\large Compte rendu Pendule de Torsion}
\end{center}

\section{Isochronisme des petites oscillations}
\begin{enumerate}
    \item Tableau des mesures ($\theta_1 > \theta_2$) 
\begin{table}[H]
    \centering
    \renewcommand{\arraystretch}{1.5}
    \begin{tabular}{|c|c|c|c|c|c|c|c|c|c|}
    \hline
    temps(s) & $t_1$ & $t_2$ & $t_3$ & $t_m$ & $\Delta t_m$ & $\Delta t_{ch}$ & $\Delta t$ & $T_0$ & $\Delta T_0$ \\
    \hline
    $\theta_1$ &  &  &  &  &  & 0,1 & &  &  \\
    \hline
    $\theta_2$ &  &  &  &  &  & 0,1 &  &  &  \\
    \hline
    \end{tabular}
\end{table}

    \item

    \item

\end{enumerate}
\section{Détermination de $C$ et $I_0$}

\begin{enumerate}
    \item Tableau des mesures ($r_1 = 5cm, r_2 = 10cm$)
\begin{table}[H]
    \centering
    \renewcommand{\arraystretch}{1.5}
    \begin{tabular}{|c|c|c|c|c|c|c|c|c|c|}
    \hline
    temps(s) & $t_1$ & $t_2$ & $t_3$ & $t_m$ & $\Delta t_m$ & $\Delta t_{ch}$ & $\Delta t$ & $T$ & $\Delta T$ \\
    \hline
    $r_1$ &  &  &  &  &  & 0,1 &  &  &  \\
    \hline
    $r_2$ &  &  &  &  &  & 0,1 &  &  &  \\
    \hline
    \end{tabular}
\end{table}

\noindent \textbf{Remarque :} Les relations utilisées étaient données en TP, à savoir :
\[ \Delta t_m = \sup|t_m - t_i| \quad , \quad T = \frac{t_m}{10} \quad , \quad \Delta T = \frac{\Delta t}{10} \quad , \quad \Delta t = \Delta t_m + \Delta t_{ch} \]
    
    \item 

    \item

\end{enumerate}

\section{Détermination de $I_{GZ}$ experimentalement}

% file: section3.tex

% ============================================================
% 1. INPUT DATA 
% ============================================================
\pgfplotstableread{
r   t1      t2      t3
3   12.45   12.55   12.50
4   12.75   12.80   12.85
5   13.15   13.20   13.25
6   13.75   13.80   13.85
7   14.45   14.50   14.55
8   15.25   15.30   15.35
9   16.90   17.00   16.95
10  22.30   22.40   22.35
}\myExperimentData

% ============================================================
% 2. CALCULATIONS
% ============================================================
\pgfplotstablecreatecol[create col/expr={(\thisrow{t1}+\thisrow{t2}+\thisrow{t3})/3}]{tm}{\myExperimentData}
\pgfplotstablecreatecol[create col/expr={max(abs(\thisrow{tm}-\thisrow{t1}), max(abs(\thisrow{tm}-\thisrow{t2}), abs(\thisrow{tm}-\thisrow{t3})))}]{Dtm}{\myExperimentData}
\pgfplotstablecreatecol[create col/set={0.1}]{const}{\myExperimentData}
\pgfplotstablecreatecol[create col/expr={0.1 + \thisrow{Dtm}}]{Dt}{\myExperimentData}
\pgfplotstablecreatecol[create col/expr={\thisrow{tm}/10}]{T}{\myExperimentData}
\pgfplotstablecreatecol[create col/expr={\thisrow{Dt}/10}]{DT}{\myExperimentData}
\pgfplotstablecreatecol[create col/expr={\thisrow{T}^2}]{Tsq}{\myExperimentData}
\pgfplotstablecreatecol[create col/expr={2 * \thisrow{T} * \thisrow{DT}}]{DTsq}{\myExperimentData}
\pgfplotstablecreatecol[create col/expr={\thisrow{r}^2}]{rsq}{\myExperimentData}
\pgfplotstablecreatecol[create col/expr={2 * \thisrow{r} * 0.1}]{Drsq}{\myExperimentData}

% ============================================================
% 3. DOCUMENT CONTENT
% ============================================================
\begin{enumerate}

    % --- (a) TABLE ---
    \item Tableau de mesures et calculs.
    \begin{table}[H]
        \centering
        \renewcommand{\arraystretch}{1.4}
        \setlength{\tabcolsep}{3.5pt}
        \resizebox{\textwidth}{!}{
            \pgfplotstabletypeset[
                columns={r, t1, t2, t3, tm, Dtm, const, Dt, T, DT, Tsq, DTsq, rsq, Drsq},
                every head row/.style={before row=\hline, after row=\hline},
                every last row/.style={after row=\hline},
                every row/.style={after row=\hline},
                columns/r/.style={column name={$r$}, precision=0, column type={|c|}},
                columns/t1/.style={column name={$t_1$}, fixed, precision=2, column type={c|}},
                columns/t2/.style={column name={$t_2$}, fixed, precision=2, column type={c|}},
                columns/t3/.style={column name={$t_3$}, fixed, precision=2, column type={c|}},
                columns/tm/.style={column name={$t_m$}, fixed, precision=2, column type={c|}},
                columns/Dtm/.style={column name={$\Delta t_m$}, fixed, precision=2, column type={c|}},
                columns/const/.style={column name={$\Delta t_{ch}$}, string type, postproc cell content/.style={@cell content={0.1}}, column type={c|}},
                columns/Dt/.style={column name={$\Delta t$}, fixed, precision=2, column type={c|}},
                columns/T/.style={column name={$T$}, fixed, precision=3, column type={c|}},
                columns/DT/.style={column name={$\Delta T$}, fixed, precision=3, column type={c|}},
                columns/Tsq/.style={column name={$T^2$}, fixed, precision=3, column type={c|}},
                columns/DTsq/.style={column name={$\Delta T^2$}, fixed, precision=3, column type={c|}},
                columns/rsq/.style={column name={$r^2$}, fixed, precision=0, column type={c|}},
                columns/Drsq/.style={column name={$\Delta r^2$}, fixed, precision=1, column type={c|}}
            ]{\myExperimentData}
        }
    \end{table}

    % --- (b) & (c) GRAPH WITH RECTANGLES ---
    \item \& c) Graphe de $T^2$ en fonction de $r^2$ avec rectangles d'incertitude.
    
    \begin{center}
        \begin{tikzpicture}
            \begin{axis}[
                width=0.85\textwidth, height=0.6\textwidth,
                title={Variation de $T^2$ en fonction de $r^2$},
                xlabel={$r^2$ ($cm^2$)}, 
                ylabel={$T^2$ ($s^2$)},
                grid=major, 
                % ADJUST THESE IF GRAPH IS SQUASHED
                xmin=0, xmax=110,
                ymin=0, ymax=6,
                legend pos=north west
            ]
            
            % 1. DRAW RECTANGLES (The complex loop)
            % This reads the table row by row and draws a box for each point
            \pgfplotstableforeachcolumnelement{rsq}\of\myExperimentData\as\xval{
                \pgfplotstablegetelem{\pgfplotstablerow}{Tsq}\of\myExperimentData \let\yval=\pgfplotsretval
                \pgfplotstablegetelem{\pgfplotstablerow}{Drsq}\of\myExperimentData \let\dx=\pgfplotsretval
                \pgfplotstablegetelem{\pgfplotstablerow}{DTsq}\of\myExperimentData \let\dy=\pgfplotsretval
                
                % Draw the rectangle: (x-dx, y-dy) to (x+dx, y+dy)
                \edef\drawrect{\noexpand\draw[blue!50, fill=blue!10, thin] 
                    (axis cs: \xval-\dx, \yval-\dy) rectangle (axis cs: \xval+\dx, \yval+\dy);}
                \drawrect
            }

            % 2. Plot the center points
            \addplot[only marks, mark=+, thick, blue] table [x=rsq, y=Tsq] {\myExperimentData};
            \addlegendentry{Points expérimentaux}

            % 3. Linear Regression (Red Line)
            \addplot[red, thick, domain=0:110] table [
                y={create col/linear regression={y=Tsq, x=rsq}}
            ] {\myExperimentData};
            \addlegendentry{Droite de régression}

            % Save slope/intercept
            \xdef\slope{\pgfplotstableregressiona}
            \xdef\intercept{\pgfplotstableregressionb}

            \end{axis}
        \end{tikzpicture}
    \end{center}

    % --- (d) ---
    \item Pentes limites et pente moyenne.
    
    L'équation de la droite est : $T^2 = \pgfmathprintnumber{\slope} \cdot r^2 + \pgfmathprintnumber{\intercept}$.
    
    Pente moyenne $p = \pgfmathprintnumber{\slope} \ s^2/cm^2$.
    
    Calcul de l'incertitude $\Delta p = \frac{|p_1 - p_2|}{2} = \dots$

    % --- (e) ---
    \item Masse $m$ et $\Delta m$.
    \[ m = \dots \quad (\pm \dots) \text{ g} \]

    % --- (f) ---
    \item Moment d'inertie $I_{GZ}$.
    \[ I_{GZ} = \dots \quad (\pm \dots) \text{ kg}\cdot m^2 \]

    % --- (g) ---
    \item Conclusion et comparaison.

\end{enumerate}% ============================================================
% 1. INPUT DATA (Edit values here)
% ============================================================
\pgfplotstableread{
r   t1      t2      t3
3   12.45   12.55   12.50
4   12.75   12.80   12.85
5   13.15   13.20   13.25
6   13.75   13.80   13.85
7   14.45   14.50   14.55
8   15.25   15.30   15.35
9   16.90   17.00   16.95
10  22.30   22.40   22.35
}\myExperimentData

% ============================================================
% 2. AUTOMATIC CALCULATIONS
% ============================================================
\pgfplotstablecreatecol[create col/expr={(\thisrow{t1}+\thisrow{t2}+\thisrow{t3})/3}]{tm}{\myExperimentData}
\pgfplotstablecreatecol[create col/expr={max(abs(\thisrow{tm}-\thisrow{t1}), max(abs(\thisrow{tm}-\thisrow{t2}), abs(\thisrow{tm}-\thisrow{t3})))}]{Dtm}{\myExperimentData}
\pgfplotstablecreatecol[create col/set={0.1}]{const}{\myExperimentData}
\pgfplotstablecreatecol[create col/expr={0.1 + \thisrow{Dtm}}]{Dt}{\myExperimentData}
\pgfplotstablecreatecol[create col/expr={\thisrow{tm}/10}]{T}{\myExperimentData}
\pgfplotstablecreatecol[create col/expr={\thisrow{Dt}/10}]{DT}{\myExperimentData}
\pgfplotstablecreatecol[create col/expr={\thisrow{T}^2}]{Tsq}{\myExperimentData}
\pgfplotstablecreatecol[create col/expr={2 * \thisrow{T} * \thisrow{DT}}]{DTsq}{\myExperimentData}
\pgfplotstablecreatecol[create col/expr={\thisrow{r}^2}]{rsq}{\myExperimentData}
\pgfplotstablecreatecol[create col/expr={2 * \thisrow{r} * 0.1}]{Drsq}{\myExperimentData}

% ============================================================
% 3. DOCUMENT CONTENT
% ============================================================
\section*{IV.3 - Détermination de $I_{GZ}$ expérimentalement}

Le balancier étant toujours avec les deux surcharges.

\begin{enumerate} % requires \usepackage{enumitem} in main.tex
% OR use \begin{enumerate} and manual \item[a)] if you don't have the package.

    % --- Question a) ---
    \item Prendre la valeur de 10 oscillations pour chaque distance de $r_3=3$ cm à $r_{10}=10$ cm. Remplir le tableau suivant.
    
    \begin{table}[H]
        \centering
        \renewcommand{\arraystretch}{1.5}
        \setlength{\tabcolsep}{3.5pt}
        \resizebox{\textwidth}{!}{
            % NO BLANK LINES INSIDE [] !
            \pgfplotstabletypeset[
                columns={r, t1, t2, t3, tm, Dtm, const, Dt, T, DT, Tsq, DTsq, rsq, Drsq},
                every head row/.style={before row=\hline, after row=\hline},
                every last row/.style={after row=\hline},
                every row/.style={after row=\hline},
                columns/r/.style={column name={$r$}, precision=0, column type={|c|}},
                columns/t1/.style={column name={$t_1$}, fixed, precision=2, column type={c|}},
                columns/t2/.style={column name={$t_2$}, fixed, precision=2, column type={c|}},
                columns/t3/.style={column name={$t_3$}, fixed, precision=2, column type={c|}},
                columns/tm/.style={column name={$t_m$}, fixed, precision=2, column type={c|}},
                columns/Dtm/.style={column name={$\Delta t_m$}, fixed, precision=2, column type={c|}},
                columns/const/.style={column name={$\Delta t_{ch}$}, string type, postproc cell content/.style={@cell content={0.1}}, column type={c|}},
                columns/Dt/.style={column name={$\Delta t$}, fixed, precision=2, column type={c|}},
                columns/T/.style={column name={$T$}, fixed, precision=3, column type={c|}},
                columns/DT/.style={column name={$\Delta T$}, fixed, precision=3, column type={c|}},
                columns/Tsq/.style={column name={$T^2$}, fixed, precision=3, column type={c|}},
                columns/DTsq/.style={column name={$\Delta T^2$}, fixed, precision=3, column type={c|}},
                columns/rsq/.style={column name={$r^2$}, fixed, precision=0, column type={c|}},
                columns/Drsq/.style={column name={$\Delta r^2$}, fixed, precision=1, column type={c|}}
            ]{\myExperimentData}
        }
    \end{table}

    % --- Question b) ---
    \item En précisant les échelles choisies, représenter graphiquement $T^2$ en fonction de $r^2$.
    
    \begin{center}
        \begin{tikzpicture}
            \begin{axis}[
                width=0.85\textwidth, height=0.55\textwidth,
                title={Graphe de $T^2$ en fonction de $r^2$},
                xlabel={$r^2$ ($m^2$)}, 
                ylabel={$T^2$ ($s^2$)},
                grid=major, 
                legend pos=north west,
            ]
            
            % Plot Points with Error Bars
            \addplot[
                only marks, mark=*, blue,
                error bars/.cd, y dir=both, y explicit, x dir=both, x explicit,
            ] table [x=rsq, y=Tsq, x error=Drsq, y error=DTsq] {\myExperimentData};
            \addlegendentry{Points expérimentaux}
    
            % Calculate and Plot Regression Line
            \addplot[
                red, thick, domain=0:110
            ] table [
                y={create col/linear regression={y=Tsq, x=rsq}}
            ] {\myExperimentData};
            \addlegendentry{Droite de régression}
    
            % Store slope/intercept for later text
            \xdef\slope{\pgfplotstableregressiona}
            \xdef\intercept{\pgfplotstableregressionb}
            \end{axis}
        \end{tikzpicture}
    \end{center}

    % --- Question c) ---
    \item Représenter sur le même graphique les rectangles d'incertitudes $\Delta(T^2)$ et $\Delta(r^2)$.
    
    \textit{(Note: Les rectangles sont représentés par les barres d'erreur bleues sur le graphique ci-dessus.)}

    % --- Question d) ---
    \item Tracer les droites limites et calculer leurs pentes $p_1$ et $p_2$, et en déduire la pente moyenne $p$.
    
    L'équation de la droite de régression calculée est : $T^2 = \pgfmathprintnumber{\slope} \cdot r^2 + \pgfmathprintnumber{\intercept}$.
    
    La pente moyenne est :
    \[ p = \pgfmathprintnumber{\slope} \ s^2/cm^2 \]
    
    Calcul de l'incertitude sur la pente :
    \[ \Delta p = \frac{|p_1 - p_2|}{2} = \dots \]
    
    % --- Question e) ---
    \item Déterminer alors la masse $m$ et son incertitude $\Delta m$.
    \[ m = \dots \]
    \[ m = ( \dots \pm \dots ) \text{ unit} \]

    % --- Question f) ---
    \item Calculer le moment d'inertie $I_{GZ}$ et $\Delta I_{GZ}$.
    \[ I_{GZ} = \dots \]
    \[ I_{GZ} = ( \dots \pm \dots ) \text{ unit} \]

    % --- Question g) ---
    \item Comparer avec $I_{GZ}$ trouvé précédemment. Conclure.
    
    \vspace{2cm} % Space for writing

\end{enumerate}

\end{document}